\documentclass[10pt,a4paper]{amsart}
\usepackage[margin=19mm]{geometry}
\usepackage{amsmath,amssymb,mathtools}
\usepackage[scr=rsfs,cal=euler]{mathalfa}
\usepackage{xfrac}
\usepackage[unicode,hidelinks,bookmarks=false]{hyperref}
\newcommand{\ie}{i.e.\ }
\newcommand{\cf}{cf.\ }
\newcommand{\resp}{respectively\ }
\newcommand{\Subsection}{\subsection}
\providecommand{\nobreakdash}{\nobreak\hbox{-}}
\setlength{\emergencystretch}{2em}
\multlinegap0pt
\usepackage{bm}
\usepackage{bbm}
\usepackage{dsfont}
\usepackage{xspace}
\usepackage{cancel}
\usepackage{mathtools}
\usepackage{amsthm}
\makeatletter
\@ifundefined{thm}{\newtheorem{thm}{Thm}}{}
\@ifundefined{lem}{\newtheorem{lem}[thm]{Lemma}}{}
\makeatother
\makeatletter
\newcommand{\PP}{\ensuremath{\mathbb{P}}}
\def\hol{{\mathcal{O}}}
\expandafter\ifx\csname myproof\endcsname\relax
\newenvironment{myproof}[2] {\paragraph{\it Proof of {#1} {#2}.}}{\hfill$\square$}
\fi
\makeatother
\title[Moduli spaces of threefolds on the Noether line]{Moduli spaces of threefolds on the Noether line}
\author{Stephen Coughlan, Yong Hu, Roberto Pignatelli, Tong Zhang}
\date{Source: arXiv:2409.17847v2 (CC BY 4.0)}
\begin{document}
\maketitle

\noindent\emph{Source and licence.} Moduli spaces of threefolds on the Noether line. Version arXiv:2409.17847v2 (CC BY 4.0), distributed under the Creative Commons Attribution 4.0 International licence, as recorded in the arXiv licence metadata for this exact version and shown on the abstract page https://arxiv.org/abs/2409.17847v2. Full rights notice: licenses/arxiv-2409.17847-CC-BY-4.0.txt.\par\smallskip

\noindent\textit{Editorial scope.} Counted unit: the Lem whose source label is \texttt{lem: dimension 10H} in the article (source0.tex lines 1241--1256, source label \texttt{lem: dimension 10H}), delivered together with the article's own complete original proof of it (source0.tex lines 1258--1336, one \texttt{proof} environment of 79 lines). No mathematics is written, repaired or re-derived: statement and proof are the article's own text line for line. Required context retained uncounted: eq: degree c (lines 870--872); tab: N=0 (lines 1205--1219); tab: N=1 (lines 1221--1237). Every definition, display and label of those lines is retained unchanged. Omitted from this package and not redistributed: 1--869, 873--1204, 1220--1220, 1238--1240, 1257--1257, 1337--2395. Nothing inside a retained excerpt is omitted or altered: every retained body is delivered exactly as the article's lines are written, so no omission receipt of any kind accompanies this package. Citations: the retained excerpts carry no citation command at all, so no bibliography entry of the article is transcribed and no reference is redistributed. Reference adaptation: none was needed and none was made; every reference inside the delivered unit resolves inside it, so no unresolved reference or citation is produced. Printed slips kept literal: no printed word, symbol, hypothesis or number of the counted statement or of its proof was altered. Layout and class: the article's own class is \texttt{amsart} with packages including amsthm, amsfonts, amsbsy, amssymb, amsmath, amscd, xy, hyperref, pgf, tikz-cd, tikz; the wrapper is the lane's pinned \texttt{amsart} editorial wrapper at 10pt on A4, which restates the theorem-like environments the retained text needs and adds the article's own macro definitions that the retained text needs (\texttt{\textbackslash PP}, \texttt{\textbackslash hol}), with extra wrapper packages (bm, bbm, dsfont, xspace, cancel, mathtools). The delivered numbering is the unit's own. Assets: no retained span contains a figure environment, a graphics-inclusion command, a PDF-page inclusion, a table or a TikZ picture; no image file is copied into this package, the package declares no asset dependency and no image file is redistributed with it. Licence: version arXiv:2409.17847v2 of this article is distributed under the Creative Commons Attribution 4.0 International licence, as recorded in the arXiv licence metadata for this exact version and shown on the abstract page https://arxiv.org/abs/2409.17847v2. A full rights notice is delivered at licenses/arxiv-2409.17847-CC-BY-4.0.txt. Characters: every retained excerpt is pure ASCII, so the delivered engine, XeLaTeX with its default Latin Modern text font, reports no missing character.
	\begin{equation} \label{eq: degree c}
		\deg c_{a_0, a_1, a_2} = N + \frac12 (a_0 + a_1) d + \frac 12 (a_1-a_0) e.
	\end{equation}

\begin{table}[ht]\caption{Vanishing monomials when $N = 0$} \label{tab: N=0}
	\begin{tabular}{ccc}
        \hline
		$d_0$ & monomials with vanishing coefficient & stratum\\
		\hline
		%$\left[ 1, \frac32 \right]$& smooth & - \\
		$<d$ & $x_0^{10},\ x_0^8y,\ x_0^6y^2,\ x_0^4y^3,\ x_0^2y^4$& terminal  \\
		$<\frac78d$ & $x_0^9x_1$ &$cA_1$\\
		$<\frac56d$ & $x_0^7x_1y$ &$cA_3$\\
		$<\frac34d$& $x_0^5x_1y^2$ &$cA_4$ \\
		$<\frac23d$ & $x_0^8x_1^2$ &$cD_6$\\
		$<\frac12d$& $x_0^6x_1^2y$,\ $x_0^3x_1y^3$ &$cE_8$ \\
		%&non-canonical & $x_0^7x_1^3$\\
	\end{tabular}
\end{table}

\begin{table}[ht] \caption{Vanishing monomials when $N = 1$ and $d \ge 3$ or when $N = 2$ and $d \ge 6$} \label{tab: N=1}
	\begin{tabular}{ccc}
        \hline
		$d_0$ & monomials with vanishing coefficient & stratum\\
		\hline
		%$\left[ 1, \frac32 \right]$& smooth & - \\
		$= d$ & $x_0^{10},\ x_0^8y,\ x_0^6y^2,\ x_0^4y^3$ & terminal \\
		$< d$ & $x_0^2y^4$ & terminal  \\
		$<\frac78d + \frac38N$ & $x_0^9x_1$ &$cA_1$\\
		$<\frac56d + \frac13N$ & $x_0^7x_1y$ &$cA_3$\\
		$<\frac34d + \frac14N$ & $x_0^5x_1y^2$ &$cA_4$ \\
		$<\frac23d + \frac13N$ & $x_0^8x_1^2$ &$cD_6$\\
		$<\frac12d + \frac14N$ & $x_0^6x_1^2y$ & $cE_7$ \\
		$<\frac12d$ & $x_0^3x_1y^3$ & $cE_8$ \\
		%&non-canonical & $x_0^7x_1^3$\\
	\end{tabular}
\end{table}

\begin{lem} \label{lem: dimension 10H}
	Suppose that $N \le 1$, $d \ge 3$ or that $N = 2$, $d \ge 6$. Then the vector space $H^0(D_z, 10H_{D_z} - 4NF_z)$ has dimension
    $$
    \begin{cases}
        125d + 36 + 36N, & \text{ if } d < d_0 \le \frac32d + \frac12N \text{ or } d_0=d, N=0;\\ 
		125d + 32 + 46N, & \text{ if } d_0 = d, N>0; \\
		155d - 30d_0 + 31 + 46N, & \text{ if } \frac78d + \frac38 N \le d_0 < d;\\ 
		162d - 38d_0 + 30 + 49N, & \text{ if } \frac56d + \frac13 N \le d_0 < \frac78d + \frac38 N;\\ 
		167d - 44d_0 + 29 + 51N, & \text{ if } \frac34d + \frac14 N \le d_0 < \frac56d + \frac13 N;\\ 
		170d - 48d_0 + 28 + 52N, & \text{ if } \frac23d + \frac13 N \le d_0 < \frac34d + \frac14 N;\\ 
		174d - 54d_0 + 27 + 54N, & \text{ if } \frac12d + \frac14 N \le d_0 < \frac23d + \frac13 N;\\
		176d - 58d_0 + 26 + 55N, & \text{ if } \frac12d \le d_0 < \frac12d + \frac14 N;\\
		177d - 60d_0 + 25 + 55N, & \text{ if } \frac14d + \frac14 N \le d_0 < \frac12d. 	 
    \end{cases}
    $$
\end{lem}

\begin{proof}
	We first observe that
	\begin{equation} \label{eq: decompose-10H}
		H^0(D_z,10H_{D_z}-4NF_z)=\bigoplus_{a_0+a_1+2a_2=10} S^{\deg c_{a_0,a_1,a_2}}(t_0,t_1)x_0^{a_0}x_1^{a_1}y^{a_2}.
	\end{equation}
    
    If $d < d_0 \le \frac12(3d + N)$ or if $d=d_0$ and $N=0$, then by Tables \ref{tab: N=0} and \ref{tab: N=1}, all the coefficients $c_{a_0,a_1,a_2}$ have non-negative degree. The number of monomials is 
	$$
	\sum_{a_2=0}^5 h^0 \left(\PP^1,\hol_{\PP^1}(10-2a_2)\right)=11+9+7+5+3+1=36.
	$$
    Thus
    \begin{align*}
    h^0(D_z,10H_{D_z}-4NF_z) & =  \sum_{\substack{a_0 + a_1 + 2a_2 = 10}} (1 + \deg c_{a_0,a_1,a_2}) \\
	& = 36 + \sum_{\substack{a_0 + a_1 + 2a_2 = 10}} \deg c_{a_0,a_1,a_2}.
    \end{align*}
    Now we replace $\deg c_{a_0,a_1,a_2}$ with its expression in \eqref{eq: degree c}. By symmetry,
    \[
    \sum_{\substack{a_0 + a_1 + 2a_2 = 10}} (a_1-a_0)=0,
    \] 
    and then
    \begin{align*}
        & \sum_{\substack{a_0 + a_1 + 2a_2 = 10}} \deg c_{a_0,a_1,a_2} = \sum \left( \frac{1}{2} (a_0 + a_1) d +N\right) \\
        = & \ \frac{1}{2}  \sum (a_0 + a_1)d  +36N = d \sum a_1 +36N= d \sum_{a_2=0}^5 \sum_{a_1=0}^{10-2a_2} a_1 + 36N \\
        = & \ d \left[ \binom{11}2 + \binom92+ \binom72+ \binom52+ \binom32 \right]+36N=125d+36N.
    \end{align*}
	This concludes the proof of the case when $d < d_0$ or $d=d_0$, $N=0$.
    
    If $d_0 = d$, $N > 0$ then the monomials $x_0^{10}$, $x_0^8y$, $x_0^6y^2$, $x_0^4y^3$ no longer appear in the equation of the branch divisor. Thus
	\begin{align*}
		h^0(D_z, 10H_{D_z} - 4NF_z) & = 125d + 36 + 36N -\sum_{k=0}^3(1 + \deg c_{10-2k,0,k})\\
		& = 125d + 36 + 36N - \left(4  - 10N \right) \\
		& = 125d  + 32 + 46N.
	\end{align*}

	If $\frac78d + \frac38 N \le d_0 < d$, then we lose the monomials 
    $x_0^{10}$, $x_0^8y$, $x_0^6y^2$, $x_0^4y^3$ and also $x_0^2y^4$. Thus
	\begin{align*}
		h^0(D_z, 10H_{D_z} - 4NF_z) & = 125d + 36 + 36N -\sum_{k=0}^4(1 + \deg c_{10-2k,0,k})\\
		& = 125d + 36 + 36N - \left(5 + 30(d_0-d) - 10N \right) \\
		& = 155d - 30d_0 + 31 + 46N.
	\end{align*}

	If $\frac56d + \frac13 N \le d_0 < \frac78d + \frac38 N$, then we also lose $x_0^9x_1$. Thus
	\begin{align*}
		h^0(D_z, 10H_{D_z} - 4NF_z) & = 155d - 30d_0 + 31 + 46N - (1 + \deg c_{9, 1, 0}) \\
		& = 162d - 38d_0 + 30 + 49N.
	\end{align*}
	
	If $\frac34d + \frac14 N \le d_0 < \frac56d + \frac13 N$, then we also lose $x_0^7x_1y$. Thus
	\begin{align*}
		h^0(D_z, 10H_{D_z} - 4NF_z) & = 162d - 38d_0 + 30 + 49N - (1 + \deg c_{7, 1, 1}) \\
		& = 167d - 44d_0 + 29 + 51N.
	\end{align*}
	
	If $\frac23d + \frac13 N \le d_0 < \frac34d + \frac14 N$, then we also lose $x_0^5x_1y^2$. Thus
	\begin{align*}
		h^0(D_z, 10H_{D_z} - 4NF_z) & = 167d - 44d_0 + 29 + 51N - (1 + \deg c_{5, 1, 2}) \\
		& = 170d - 48d_0 + 28 + 52N.
	\end{align*}
	
	If $\frac12d + \frac14 N \le d_0 < \frac23d + \frac13 N$, then we also lose $x_0^8x_1^2$. Thus
	\begin{align*}
		h^0(D_z, 10H_{D_z} - 4NF_z) & = 170d - 48d_0 + 28 + 52N - (1 + \deg c_{8, 2, 0}) \\
		& = 174d - 54d_0 + 27 + 54N.
	\end{align*}
	
	If $\frac12d \le d_0 < \frac12d + \frac14 N$, then we also lose $x_0^6x_1^2y$. Thus
	\begin{align*}
		h^0(D_z, 10H_{D_z} - 4NF_z) & = 174d - 54d_0 + 27 + 54N - (1 + \deg c_{6, 2, 1}) \\
		& = 176d - 58d_0 + 26 + 55N.
	\end{align*}
	
	If $\frac14(d + N) \le d_0 < \frac12d$, then we also lose $x_0^3x_1y^3$. Thus
	\begin{align*}
		h^0(D_z, 10H_{D_z} - 4NF_z) & = 176d - 58d_0 + 26 + 55N - (1 + \deg c_{3, 1, 3}) \\
		& = 177d - 60d_0 + 25 + 55N.
	\end{align*}
	This concludes the proof.
\end{proof}

\end{document}
